\documentclass[11pt,a4paper]{article}
\usepackage{xcolor}
\usepackage[margin=1in]{geometry}
\usepackage{amsmath,amssymb,amsthm,mathtools}
\usepackage{bm}
\usepackage{booktabs}
\usepackage{array}
\usepackage{enumitem}
\usepackage{xcolor}
\usepackage{hyperref}
\usepackage{tikz}
\usepackage{listings}

\hypersetup{
    colorlinks=true,
    linkcolor=blue!60!black,
    citecolor=blue!60!black,
    urlcolor=blue!60!black
}

\newcommand{\F}{\mathbb{F}}
\newcommand{\wt}{\operatorname{wt}}
\newcommand{\Image}{\operatorname{Im}}
\newcommand{\Ber}{\mathrm{Ber}}
\newcommand{\U}{\mathcal{U}}
\newcommand{\E}{\mathbb{E}}
\newcommand{\Prob}{\Pr}
\newcommand{\Z}{\mathbb{Z}}

\title{\textbf{Understanding the d-LIN Problem}\\
\large Sparse Noisy Linear Equations over $\mathbb{F}_2$}
\author{Prajjwal Saxena (with the help of GPT 5.6 Luna)}
\date{}

\begin{document}

\maketitle



\tableofcontents
\newpage

% ============================================================
\section{The Basic Idea}

The central d-LIN equation is\begin{equation}
    \boxed{b = Mx + e}
    \label{eq:main}
\end{equation}

where the arithmetic is performed over the binary field \[
    \F_2 = \{0,1\}.
\]

The objects are:
\begin{itemize}
    \item $M$: a random $d$-sparse matrix;
    \item $x$: a hidden binary assignment;
    \item $e$: a random noise vector;
    \item $b$: the observed right-hand-side vector.
\end{itemize}

The essential intuition is:
\textbf{\[
\text{d-LIN} =
\text{sparse linear equations}
+
\text{noise}
\quad\text{(over) }\F_2.
\]}

Without the noise $e$, the problem is ordinary linear algebra and can
be solved efficiently. 

% ============================================================

\section{The Binary Field \texorpdfstring{$\F_2$}{F2}}
The field $\F_2$ contains only
\[
    \F_2 = \{0,1\}.
\]
Addition is modulo $2$:\begin{align}
0+0 &= 0,\\
0+1 &= 1,\\
1+0 &= 1,\\
1+1 &= 0.
\end{align}

Thus addition in $\F_2$ is equivalent to XOR: \begin{equation}
    a+b \pmod 2 = a\oplus b.
\end{equation}

Subtraction is also the same as addition:\begin{equation}
    a-b = a+b \pmod 2.
\end{equation}

Consequently, a linear equation such as\begin{equation}
    x_1+x_4+x_7 = 1
\end{equation}

can equivalently be written as\begin{equation}
    x_1\oplus x_4\oplus x_7 = 1.
\end{equation}
This makes d-LIN a natural binary parity or XOR constraint problem.

% ============================================================
\section{What Is a d-LIN Instance?}
An instance with \textbf{$m$ clauses} and $n$ \textbf{variables} consists of\begin{equation}
    M\in\F_2^{m\times n}
\end{equation}
and
\begin{equation}
    b\in\F_2^m.
\end{equation}
The unknown assignment is
\begin{equation}
    x=
    \begin{pmatrix}
    x_1\\
    x_2\\
    \vdots\\
    x_n
    \end{pmatrix}
    \in\F_2^n.
\end{equation}
The $m$ rows of $M$ correspond to the $m$ clauses.

% ------------------------------------------------------------


\subsection{The Meaning of \texorpdfstring{$m$}{m}}
The integer $m$ is the number of equations or clauses.

For example, $m=4$ means that there are four constraints:\begin{align}
C_1 &: \quad \text{first equation},\\
C_2 &: \quad \text{second equation},\\
C_3 &: \quad \text{third equation},\\
C_4 &: \quad \text{fourth equation}.
\end{align}

% ------------------------------------------------------------


\subsection{The Meaning of \texorpdfstring{$n$}{n}}

The integer $n$ is the number of binary variables:

\begin{equation}
    x_1,x_2,\ldots,x_n\in\{0,1\}.
\end{equation}

There are therefore $2^{n}$ possible assignments.


% ============================================================


\section{What Does d-Sparse Mean?}

A matrix $M$ is called $d$-sparse in the ABW09 setting when every
row contains exactly $d$ ones.

For example, a $3$-sparse matrix can be

\begin{equation}
M =
\begin{pmatrix}
1&1&0&1&0\\
0&1&1&0&1\\
1&0&1&1&0\\
0&1&0&1&1
\end{pmatrix}.
\end{equation}

Every row contains exactly three $1$s.

The first row,

\[
(1,1,0,1,0),
\]

represents the equation

\begin{equation}
    x_1+x_2+x_4=b_1.
\end{equation}

The second row represents

\begin{equation}
    x_2+x_3+x_5=b_2.
\end{equation}

Thus, for $d=3$,

\begin{equation}
    \text{3-LIN is a collection of 3-variable XOR equations.}
\end{equation}

The zero entries mean that the corresponding variable does not
participate in that clause.

% ============================================================
\section{Matrix Form of the Clauses}

Let

\begin{equation}
M =
\begin{pmatrix}
M_{11}&M_{12}&\cdots&M_{1n}\\
M_{21}&M_{22}&\cdots&M_{2n}\\
\vdots&\vdots&\ddots&\vdots\\
M_{m1}&M_{m2}&\cdots&M_{mn}
\end{pmatrix}.
\end{equation}

Then

\begin{equation}
Mx =
\begin{pmatrix}
\sum_{j=1}^n M_{1j}x_j\\
\sum_{j=1}^n M_{2j}x_j\\
\vdots\\
\sum_{j=1}^n M_{mj}x_j
\end{pmatrix}.
\end{equation}

Each row gives one clause.

For example,

\begin{equation}
M_1=(1,1,0,1,0)
\end{equation}

gives

\begin{equation}
M_1x=x_1+x_2+x_4.
\end{equation}

Therefore

\begin{equation}
Mx=b
\end{equation}

is simply a compact representation of all $m$ equations.

% ============================================================
\section{Exact Satisfiability}

If \begin{equation}
    Mx=b,
\end{equation}
then every clause is satisfied.
Equivalently, \begin{equation}
    Mx-b=0.
\end{equation}

Hence

\begin{equation}
    \wt(Mx-b)=0.
\end{equation}

This is an exact d-LIN solution.

Exact linear systems over $\F_2$ can be solved efficiently
using Gaussian elimination. 

% ============================================================
\section{Hamming Weight and Violated Clauses}

For a binary vector \begin{equation}
v=(v_1,\ldots,v_m),
\end{equation}

its Hamming weight is \begin{equation}
    \wt(v)=\left|\{i:v_i=1\}\right|.
\end{equation}

For example,\begin{equation}
v=(0,1,0,1,1)
\end{equation}
has\begin{equation}
    \wt(v)=3.
\end{equation}
Now consider\begin{equation}
    Mx-b.
\end{equation}
For the $i$th clause, \begin{equation}
    (Mx-b)_i=0
\end{equation}

means that the clause is satisfied, while

\begin{equation}
    (Mx-b)_i=1
\end{equation} means that it is violated.

Therefore, \textbf{\begin{equation}
    \boxed{
    \wt(Mx-b)
    =
    \text{number of violated clauses}.
    }
\end{equation}}

% ============================================================
\section{Almost Satisfiable d-LIN}

An instance is called $(1-\epsilon)$-satisfiable if there exists an
assignment $x$ such that

\begin{equation}
    \wt(Mx-b)\leq \epsilon m.
    \label{eq:satisfiable}
\end{equation}

There are $m$ total clauses and at most $\epsilon m$ violated clauses.
Hence at least

\begin{equation}
    m-\epsilon m
    =
    (1-\epsilon)m
\end{equation}

clauses are satisfied.

The fraction satisfied is therefore at least

\begin{equation}
    \frac{(1-\epsilon)m}{m}
    =
    1-\epsilon.
\end{equation}

Thus,

\begin{equation}
    \boxed{
    \wt(Mx-b)\leq\epsilon m
    \iff
    \text{at least }(1-\epsilon)\text{ of the clauses are satisfied}.
    }
\end{equation}

For example, if $m=1000$ and $\epsilon=0.01$, then at most 10
clauses are violated, and at least 990 are satisfied.

% ============================================================
\section{Bernoulli Noise}

ABW09 considers a Bernoulli noise distribution denoted by : 
    $\Ber_\epsilon^m$

An error vector \begin{equation}
    e=(e_1,\ldots,e_m)
\end{equation}

is sampled such that \begin{equation}
    \Prob(e_i=1)=\epsilon
\end{equation}
and
\begin{equation}
    \Prob(e_i=0)=1-\epsilon.
\end{equation}

The bits are sampled independently.
Thus, \begin{equation}
    e_i\sim\Ber(\epsilon)
\end{equation}

independently for every $i$.

The expected number of errors is

\begin{equation}
    \E[\wt(e)]
    =
    \sum_{i=1}^m\E[e_i]
    =
    m\epsilon.
\end{equation}

Therefore the expected fraction of corrupted equations is $\epsilon$.

% ============================================================
\section{Generating an Almost-Satisfiable Instance}

A natural generation procedure is:

\begin{enumerate}
    \item Choose a random $d$-sparse matrix
    \[
        M\leftarrow M_{m,n,d}.
    \]

    \item Choose a uniformly random binary assignment
    \[
        x\leftarrow \U_n.
    \]

    \item Choose independent Bernoulli noise
    \[
        e\leftarrow\Ber_\epsilon^m.
    \]

    \item Publish
    \[
        \boxed{b=Mx+e}.
    \]
\end{enumerate}

All operations are over $\F_2$.

Without noise, $Mx$ belongs exactly to the set of valid right-hand
sides generated by $M$. The noise moves it away from that set.

% ============================================================
\section{The Image of a Matrix}

The image of $M$ is

\begin{equation}
    \Image(M)
    =
    \{Mx:x\in\F_2^n\}.
\end{equation}

Thus $\Image(M)$ is the set of all vectors that can be produced by
multiplying $M$ by some binary vector $x$.

For every assignment $x$,

\begin{equation}
    Mx\in\Image(M).
\end{equation}

Since

\begin{equation}
    b=Mx+e,
\end{equation}

the observed vector $b$ is obtained by taking one point in
$\Image(M)$ and adding noise.

% ------------------------------------------------------------
\subsection{Distance from the Image}

The Hamming distance between two binary vectors is

\begin{equation}
    d_H(u,v)
    =
    \wt(u-v).
\end{equation}

The distance of $b$ from the image of $M$ is

\begin{equation}
    d_H(b,\Image(M))
    =
    \min_{y\in\Image(M)}d_H(b,y).
\end{equation}

The construction

\[
b=Mx+e
\]

guarantees that $b$ is at distance at most $\wt(e)$ from the
particular image point $Mx$:

\begin{equation}
    d_H(b,Mx)=\wt(e).
\end{equation}

Since

\[
\E[\wt(e)]=\epsilon m,
\]

the expected absolute distance is $\epsilon m$, or normalized by $m$,

\begin{equation}
    \boxed{
    \text{expected relative distance}\approx\epsilon.
    }
\end{equation}

% ============================================================
\section{A Complete Numerical Example}
\begin{equation}
M=
\begin{pmatrix}
1&1&0&1&0\\
0&1&1&0&1\\
1&0&1&1&0\\
0&1&0&1&1
\end{pmatrix},
\qquad
b=
\begin{pmatrix}
1\\1\\0\\1
\end{pmatrix}.
\end{equation}
Choose
\begin{equation}
x=
\begin{pmatrix}
1\\0\\1\\0\\1
\end{pmatrix}.
\end{equation}
Then
\begin{equation}
Mx=
\begin{pmatrix}
1\\0\\0\\1
\end{pmatrix}.
\end{equation}

Suppose

\begin{equation}
e=
\begin{pmatrix}
0\\1\\0\\0
\end{pmatrix}.
\end{equation}

Then

\begin{equation}
b=Mx+e
=
\begin{pmatrix}
1\\0\\0\\1
\end{pmatrix}
+
\begin{pmatrix}
0\\1\\0\\0
\end{pmatrix}
=
\begin{pmatrix}
1\\1\\0\\1
\end{pmatrix}.
\end{equation}

Thus

\begin{equation}
Mx-b
=
\begin{pmatrix}
0\\1\\0\\0
\end{pmatrix},
\end{equation}

and

\begin{equation}
\wt(Mx-b)=1.
\end{equation}

There are four clauses, so the assignment satisfies

\begin{equation}
    1-\frac14=\frac34
\end{equation}

of the clauses.

Hence the instance is at least $3/4$-satisfiable.

% ============================================================
\section{Where Does the Hardness Come From?}

It is important to distinguish three problems.

\subsection{Exact linear system}

\begin{equation}
    Mx=b.
\end{equation}

This can be solved efficiently by Gaussian elimination over $\F_2$.

Therefore:

\begin{equation}
    \boxed{\text{linear algebra itself is not hard}.}
\end{equation}

\subsection{Noisy system}

The cryptographic problem is instead

\begin{equation}
    b=Mx+e.
\end{equation}

The attacker receives $(M,b)$ but not $x$ or $e$.

The attacker must find an assignment that explains most of the
observed equations.

Equivalently, the task can be written approximately as\begin{equation}
    \boxed{
    \text{find }x\text{ such that }
    d_H(Mx,b)\leq\epsilon m.
    }
\end{equation}
This is a noisy constraint-satisfaction problem.
\subsection{Unknown locations of the errors}
The attacker does not know which equations are corrupted.

For example, if
\[
m=1000,\qquad \epsilon=0.1,
\]
then approximately 100 equations are noisy, but their positions are
unknown.

The attacker must simultaneously infer the hidden assignment and
cope with the unknown corrupted constraints.

\subsection{Exponential assignment space}
There are\begin{equation}
    2^n
\end{equation}
possible binary assignments.
Brute-force search therefore takes exponential time in $n$.

The cryptographic assumption is that no known efficient algorithm can
recover the hidden assignment, or otherwise exploit the generated
instance, for the selected parameter regime.

% ============================================================
\section{What Exactly Is Assumed to Be Hard?}
The statement should be phrased carefully.

The assumption is not
\begin{equation}
    \text{d-LIN is mathematically proven to be hard}.
\end{equation}
Rather, it is a cryptographic hardness assumption of the form
\begin{equation}
\boxed{
\begin{minipage}{0.85\textwidth}
For appropriately chosen parameters, a random noisy sparse system
\[
b=Mx+e
\]
over $\F_2$ cannot be efficiently solved by a probabilistic
polynomial-time algorithm.
\end{minipage}}
\end{equation}
This is analogous in spirit to the way cryptography relies on
assumptions such as LPN, LWE, and related problems.

\textbf{ABW09 provides evidence against several classes of attacks, including \textcolor{red}{myopic attacks, low-degree polynomial approaches, $AC^0$ algorithms,and Lasserre/SDP-type methods}. Such evidence is not a proof that
d-LIN is outside $\mathbf P$.}

% ============================================================
\section{Relation to LPN}

Learning Parity with Noise (LPN) has the general form

\begin{equation}
    b=Ax+e
    \pmod 2.
\end{equation}

This looks extremely similar to d-LIN.

The main structural difference is that in d-LIN the matrix is
explicitly sparse:\begin{equation}
    \wt(M_i)=d
\end{equation} for every row $M_i$.

Thus, d-LIN can be viewed as a sparse noisy parity-learning regime.

For $d=3$,

\begin{equation}
    x_i+x_j+x_k=b_\ell
\end{equation}

is a noisy 3-XOR constraint.

The common structure is

\begin{equation}
    \boxed{\text{hidden linear relation}+\text{noise}.}
\end{equation}

The exact assumptions, parameter regimes, and known algorithms are
nevertheless different.

% ============================================================
\section{Relation to LWE}

The Learning With Errors problem has the form

\begin{equation}
    b=As+e\pmod q,
\end{equation}

where typically

\begin{equation}
    A\in\Z_q^{m\times n}.
\end{equation}

Compare:

\begin{center}
\begin{tabular}{lll}
\toprule
 & d-LIN & LWE\\
\midrule
Domain & $\F_2$ & $\Z_q$\\
Matrix & sparse & generally random/dense\\
Secret & binary $x$ & $s\in\Z_q^n$\\
Noise & Bernoulli & small modular/integer noise\\
Equation & $b=Mx+e$ & $b=As+e\pmod q$\\
\bottomrule
\end{tabular}
\end{center}

The noiseless versions of both problems are easy linear algebra.
The difficulty comes from the noisy version.

However, d-LIN and LWE are not simply the same problem with different
notation.

% ============================================================
\section{Relation to MLWE}

Module-LWE introduces a module or polynomial-ring structure.

A schematic MLWE equation is

\begin{equation}
    b=As+e
\end{equation}

over a ring such as

\begin{equation}
    R_q=\Z_q[X]/(f(X)).
\end{equation}

Thus:

\begin{center}
\begin{tabular}{lll}
\toprule
Problem & Algebraic setting & Main structure\\
\midrule
d-LIN & $\F_2$ & sparse binary matrix\\
LPN & $\F_2$ & noisy parity learning\\
LWE & $\Z_q$ & modular linear equations\\
MLWE & polynomial/module ring & structured modular equations\\
\bottomrule
\end{tabular}
\end{center}

All share the high-level pattern

\begin{equation}
    \boxed{\text{recover hidden structure from noisy linear information}.}
\end{equation}

But their mathematical hardness assumptions are different.

% ============================================================
\section[Why Does ABW09 Use F2?]{Why Does ABW09 Use $\F_2$?}

There is nothing mathematically inevitable about $\F_2$.
One could define sparse noisy linear equations over a larger finite
field $\F_q$.

ABW09 chooses $\F_2$ for structural and conceptual reasons.

\subsection{Boolean variables}

Over $\F_2$,

\begin{equation}
    x_i\in\{0,1\}.
\end{equation}

Thus each equation is a Boolean parity constraint.

For example,

\begin{equation}
    x_1+x_4+x_7=1
\end{equation}

is

\begin{equation}
    x_1\oplus x_4\oplus x_7=1.
\end{equation}

This makes the problem naturally combinatorial.

\subsection{Clean interpretation of noise}

A Bernoulli error bit has only two possibilities:

\begin{equation}
e_i=
\begin{cases}
0,&\text{equation unchanged},\\
1,&\text{equation right-hand side flipped}.
\end{cases}
\end{equation}

Thus

\begin{equation}
    b=Mx+e
\end{equation}

literally means that a small fraction of the parity equations have
their output bit flipped.

\subsection{Cancellation}

The characteristic-two identity

\begin{equation}
    x+x=0
\end{equation}

gives especially clean cancellation.

Suppose a subset $S$ of rows satisfies

\begin{equation}
    \sum_{i\in S}M_i=0.
\end{equation}

Then for

\begin{equation}
    b=Mx+e
\end{equation}

we obtain

\begin{align}
\sum_{i\in S}b_i
&=
\sum_{i\in S}(M_ix+e_i)\\
&=
\left(\sum_{i\in S}M_i\right)x
+
\sum_{i\in S}e_i\\
&=
\sum_{i\in S}e_i.
\end{align}

The unknown $x$ cancels completely.

This cancellation property is central to the hidden linear dependency
used as a trapdoor in the ABW public-key construction.

% ============================================================
\section{Could We Use \texorpdfstring{$\F_q$}{Fq} Instead?}

In principle, yes.

One could define

\begin{equation}
    b=Mx+e\pmod q,
\end{equation}

with

\begin{equation}
    M\in\F_q^{m\times n}
\end{equation}

and sparse rows.

However, this would define a different problem.

Over $\F_2$:

\begin{equation}
    x_i+x_j+x_k=b
\end{equation}

is a Boolean XOR constraint.

Over $\F_q$:

\begin{equation}
    x_i+x_j+x_k=b\pmod q
\end{equation}

is a constraint over $q$ possible values per variable.

The reductions, combinatorial interpretation, noise distribution,
parameter choices, and attack landscape can therefore change.

Consequently, replacing $\F_2$ by $\F_q$ does not automatically produce
``the same d-LIN problem but stronger.'' It creates a related but
distinct hardness assumption.

% ============================================================
\section{Connection to the ABW09 Public-Key Construction}

The sparse noisy linear assumption is not introduced merely as an
abstract puzzle. ABW09 uses it to construct public-key encryption.

The central relation remains

\begin{equation}
    b=Mx+e.
\end{equation}

The construction exploits a secret short linear dependency among
rows of the public matrix:

\begin{equation}
    \sum_{i\in S}M_i=0.
\end{equation}

When the corresponding ciphertext components are summed, the
unknown term involving $x$ disappears:

\begin{equation}
    \sum_{i\in S}b_i
    =
    \sum_{i\in S}e_i.
\end{equation}

The legitimate decryptor knows the hidden dependency $S$, while an
attacker is not supposed to be able to exploit the public sparse
system in the same way.

This is the connection between the abstract d-LIN assumption and the
cryptographic trapdoor construction.

% ============================================================
\section{Search and Decision Perspectives}

A useful distinction is between a \emph{search} problem and a
\emph{decision} problem.
\subsection{Search} Given \begin{equation}
    (M,b)
\end{equation}
generated according to \begin{equation}
    b=Mx+e,
\end{equation} recover the hidden $x$ or otherwise find an assignment explaining
most of the equations.

\subsection{Decision}

A decision formulation asks whether a given sample is distributed
like a valid noisy d-LIN instance or instead comes from some
appropriate reference distribution.

ABW09 develops reductions connecting search, approximate search,
prediction, and distributional formulations in order to support the
cryptographic construction.


% ============================================================
\section{What Makes d-LIN Hard? A Careful Statement}

The safest conceptual statement is:
\begin{equation}
\boxed{
\parbox{0.9\textwidth}{
\centering % Optional: centers the text inside the box
The assumed hardness comes from recovering a hidden assignment
from random sparse equations whose outputs have been independently
corrupted by noise.
}}
\end{equation}


It is \emph{not} correct to say that ordinary linear algebra over
$\F_2$ is hard.

It is also not correct to say that ABW09 proves that d-LIN cannot be
solved efficiently.

The paper provides cryptographic assumptions and evidence against
specific classes of algorithms. As with other cryptographic
assumptions, security depends on the parameter regime and on the
absence of efficient attacks.

% ============================================================
\section{A Useful Mental Model}

The progression can be remembered as follows: \[
\boxed{
b=Mx
}
\] means:\[
\text{exact linear equations}
\quad\Longrightarrow\quad
\text{easy by Gaussian elimination}.
\]Add noise: \[
\boxed{
b=Mx+e
}
\] and now:
\[
\text{which equations are wrong?}
\]
becomes part of the problem. 
Make $M$ sparse:
\[
\boxed{
\operatorname{rowwt}(M)=d
}
\]
and the equations become local parity constraints:
\[
x_{i_1}\oplus x_{i_2}\oplus\cdots\oplus x_{i_d}=b_j.
\]
Therefore the central picture is
\[
\boxed{
\text{random sparse parity constraints}
+
\text{random noise}
+
\text{hidden assignment}.
}
\]

% ============================================================
\section{Summary Table}

\begin{center}
\begin{tabular}{p{0.22\textwidth}p{0.68\textwidth}}
\toprule
Symbol / Term & Meaning\\
\midrule
$m$ & Number of clauses/equations\\
$n$ & Number of binary variables\\
$d$ & Number of ones in every row of $M$\\
$M$ & $m\times n$ sparse binary constraint matrix\\
$x$ & Hidden $n$-bit assignment\\
$b$ & Observed $m$-bit right-hand-side vector\\
$e$ & Noise/error vector\\
$\Ber_\epsilon^m$ & $m$ independent Bernoulli bits with $\Prob(e_i=1)=\epsilon$\\
$\wt(v)$ & Hamming weight, i.e. number of ones in $v$\\
$\Image(M)$ & Set $\{Mx:x\in\F_2^n\}$\\
$\epsilon m$ & Expected number of corrupted clauses\\
$1-\epsilon$ & Expected/guaranteed satisfiable fraction in the corresponding bound\\
\bottomrule
\end{tabular}
\end{center}

% ============================================================
\section{Key Equations to Remember}

The most important equations are:

\begin{align}
    &\boxed{M\in\F_2^{m\times n}}
    &&\text{binary constraint matrix},\\[4pt]
    &\boxed{\wt(M_i)=d}
    &&\text{$d$ ones per row},\\[4pt]
    &\boxed{b=Mx+e}
    &&\text{noisy d-LIN instance},\\[4pt]
    &\boxed{\wt(Mx-b)}
    &&\text{number of violated clauses},\\[4pt]
    &\boxed{\wt(Mx-b)\leq\epsilon m}
    &&\text{at least }1-\epsilon\text{ clauses satisfied},\\[4pt]
    &\boxed{\Image(M)=\{Mx:x\in\F_2^n\}}
    &&\text{all noiseless outputs},\\[4pt]
    &\boxed{\E[\wt(e)]=\epsilon m}
    &&\text{expected number of errors}.
\end{align}

The broader learning-with-noise comparison is:

\begin{align}
    \text{d-LIN:}\quad
    &b=Mx+e\pmod2,
    &&M\text{ sparse},\\
    \text{LPN:}\quad
    &b=Ax+e\pmod2,
    &&A\text{ generally random binary},\\
    \text{LWE:}\quad
    &b=As+e\pmod q,
    &&A\in\Z_q^{m\times n},\\
    \text{MLWE:}\quad
    &b=As+e,
    &&A,s,e\text{ over a module/ring}.
\end{align}

The common high-level idea is:

\[
\boxed{
\text{hidden linear structure}
+
\text{noise}
\longrightarrow
\text{computational hardness assumption}.
}
\]

\end{document}
